\section{The AI assistant}
\label{sec:ai}

The assistant is built around one idea: it should change your document, not
describe a change and leave you to transcribe it. Everything below follows from
that.

\subsection{Setting it up}

Open \ui{Settings}, choose a provider, paste its key, and press
\ui{Test Connection}. Supported providers are OpenAI, Anthropic, DeepSeek,
OpenRouter, Grok, OpenCode Go, and Ollama --- the last running locally and
needing no key at all, which is the option to reach for if you would rather
your draft never left the machine.

Keys are stored in \texttt{\textasciitilde/.latex-editor/settings.json}.

\subsection{Model and speed}
\label{sec:ai-model}

Skills work on short passages, so what matters is latency, not raw capability.
The thing that dominates latency is whether the model \emph{deliberates} before
answering --- and most current models do, at length.

Proofreading one section of this guide, identical text and prompt each time:

\begin{table}[htbp]
  \centering
  \caption{Measured on this document, via OpenCode Go}
  \label{tab:models}
  \begin{tabular}{lrr}
    \toprule
    \textbf{Configuration} & \textbf{Time} & \textbf{Thinking} \\
    \midrule
    \texttt{deepseek-v4-pro}                        & 150\,s & 40{,}700 chars \\
    \texttt{deepseek-v4-flash}                      & 45\,s  & 22{,}600 chars \\
    \texttt{deepseek-v4-flash}, deliberation off    & 6\,s   & 0 \\
    \bottomrule
  \end{tabular}
\end{table}

Twenty-five times faster end to end, for an edit of the same quality. The two
settings that get you there are both on by default:

\begin{itemize}[leftmargin=*]
  \item \textbf{A fast model.} Reach for a heavier one only where the reasoning
        earns its keep, which proofreading a paragraph does not.
  \item \textbf{\ui{Answer directly, without deliberating}}, in
        \ui{Settings} $\rightarrow$ \ui{Speed}. Turn it off when you want the
        model to think a problem through.
\end{itemize}

If a model does deliberate, the panel streams its chain of thought under
\ui{Thinking} so you can see it working rather than hung. That text is never
part of the answer and is never written into your document.

\paragraph{A note on Max Tokens.}
Deliberation is charged against the same output budget as the answer. A model
that thinks for twelve thousand tokens on a four thousand token budget produces
nothing at all --- so the default is 16{,}384, and the editor tells you plainly
when a run ended that way rather than showing an empty result.

\subsection{The inline assistant}
\label{sec:inline}

Most edits should never involve the sidebar at all. Put the cursor in a
paragraph and press \kbd{Cmd+K}.

A bar opens under the text with the passage it is about to change tinted behind
it. You did not have to select anything: with no selection, \kbd{Cmd+K} takes
the paragraph you are standing in, bounded by blank lines and by
\verb|\begin|/\verb|\end|, so pressing it inside a table body works on the row
rather than swallowing the environment. Select first if you want a different
scope.

Then either type what you want --- \emph{say this in plainer language},
\emph{cut the hedging} --- or click one of the presets. The answer streams in,
you get a diff, and:

\begin{itemize}[leftmargin=*]
  \item \kbd{Cmd+Enter} accepts it.
  \item \kbd{Esc} discards it, or stops a run in progress. It works wherever
        the keyboard focus happens to be, including after you have clicked back
        into the document --- and there is a close button on the bar for the
        same purpose.
  \item \kbd{Enter} refines: type \emph{shorter}, \emph{less formal}, and it
        runs again --- against the \emph{original} text, not its own last
        answer, so refinements do not compound edits on top of edits.
\end{itemize}

Nothing is written until you accept, and \kbd{Cmd+Z} undoes an accept like any
other edit.

\subsection{Reading a diff}

Where a line was rewritten rather than replaced wholesale, only the words that
actually changed are tinted --- the rest of the line is dimmed:

\begin{lstlisting}
- are all coloured differently, [\verb|\begin|/\verb|\end|] blocks fold, and
+ are all coloured differently, [code] blocks fold, and
\end{lstlisting}

This matters more than it sounds. A plain line diff shows a one-word fix as a
whole line deleted and a whole line added, and leaves you to find the
difference yourself. Pairing only happens where it helps: if the two versions
share no wording, or the changes are scattered across the whole line, you get
the two plain lines instead, because a word diff there is confetti.

\subsection{The loop}

\begin{enumerate}[leftmargin=*]
  \item \textbf{Select} the passage you want changed.
  \item Open \ui{AI} in the sidebar and pick a skill.
  \item Press \ui{Run}. The answer streams in; the button becomes
        \ui{Cancel} with a running timer.
  \item \textbf{Read the diff.} Removed lines in red, added in green,
        unchanged stretches collapsed to \texttt{$\cdots$ n unchanged lines}.
  \item Press \ui{Apply} --- or do not.
\end{enumerate}

\subsection{Selection is the point}

Before you press \ui{Run}, the panel states what it is about to send:
\emph{selection (4 lines)}, or \emph{whole file (312 lines)}. That line is the
one to check.

With a paragraph selected, only that paragraph goes to the model, and
\ui{Apply} writes the result back over exactly those lines --- the rest of your
document is untouched, and you pay for a paragraph instead of a thesis. With
nothing selected the whole file is sent, which is right for
\ui{Generate Abstract} and wasteful for \ui{Proofread Section}.

The range is captured when you press \ui{Run}, not when you press \ui{Apply}.
You are free to scroll and click around while the model works. And if you
switch to a different file, \ui{Apply} is disabled with a note explaining why
--- rather than writing your revised paragraph into whatever document happens
to be open.

\subsection{Skills do different things}

Skills are not interchangeable, and the panel treats them differently.

\begin{table}[htbp]
  \centering
  \caption{Skill kinds and what the panel offers for each}
  \label{tab:skills}
  \begin{tabular}{lll}
    \toprule
    \textbf{Kind} & \textbf{Skills} & \textbf{You get} \\
    \midrule
    Rewrite & Proofread, Tighten, Improve Equation,   & a diff, and \ui{Apply} \\
            & Fix Errors, Fix Bibliography            & \\[2pt]
    Generate & Generate Table, Paste Data             & the \LaTeX, and \\
             & $\rightarrow$ Table, Generate Abstract, & \ui{Insert at cursor} \\
             & Executive Summary, Notes               & \\
             & $\rightarrow$ Prose, Plain Text
               $\rightarrow$ \LaTeX                   & \\[2pt]
    Explain & Explain Equation, Summarize Section,    & prose, and \ui{Copy} \\
            & Consistency Check,                      & only \\
            & Suggest Paper Structure                 & \\
    \bottomrule
  \end{tabular}
\end{table}

There is no \ui{Apply} button on \ui{Explain Equation}, because an explanation
is for you to read, not to paste into your paper. A diff would be meaningless
there too, so none is shown.

\subsection{It reads your project}

Each skill declares what it needs to know, and the panel gathers it before
sending --- listing what it collected under \emph{Also sends}, so the cost is
never hidden from you.

\begin{table}[htbp]
  \centering
  \caption{Context attached to each skill}
  \label{tab:context}
  \begin{tabular}{ll}
    \toprule
    \textbf{Skill} & \textbf{Also sends} \\
    \midrule
    Proofread Section, Tighten     & labels, citation keys \\
    Inline edit (\kbd{Cmd+K})      & preamble, labels, citation keys \\
    Improve Equation               & preamble \\
    Generate Table, Paste Data
      $\rightarrow$ Table          & preamble \\
    Notes $\rightarrow$ Prose,
      Plain Text $\rightarrow$ \LaTeX & preamble, labels, citation keys \\
    Fix Compilation Errors         & the errors, preamble \\
    Fix Bibliography               & citation keys \\
    Consistency Check              & labels \\
    Suggest Paper Structure        & project file tree \\
    Explain, Summarize,
      Executive Summary            & nothing \\
    \bottomrule
  \end{tabular}
\end{table}

This is what stops the two most irritating failure modes. A model that has been
shown your preamble will not hand you a \verb|\toprule| for a document that
never loaded \texttt{booktabs}. A model that has been given the keys in your
\texttt{.bib} will cite \texttt{knuth1984} rather than inventing
\texttt{smith2019}.

\ui{Explain Equation} sends nothing extra, because nothing extra would help ---
paying for a file tree to explain a formula is waste.

\subsection{Worked example: fixing a compilation error}

This is the skill where the difference is starkest, because the editor already
has everything needed.

Suppose you typed:

\begin{lstlisting}
\includegrpahics[width=0.8\textwidth]{plot.pdf}
\end{lstlisting}

Press \ui{Typeset} and \ui{Issues} reports
\emph{Undefined control sequence}. Now:

\begin{enumerate}[leftmargin=*]
  \item Select the offending line.
  \item \ui{AI} $\rightarrow$ \ui{Fix Compilation Errors} $\rightarrow$ \ui{Run}.
\end{enumerate}

You are not asked to paste the log. The panel shows
\emph{Also sends: 1 error, preamble}, having taken the error straight from the
run you just did. The diff comes back:

\begin{lstlisting}
- \includegrpahics[width=0.8\textwidth]{plot.pdf}
+ \includegraphics[width=0.8\textwidth]{plot.pdf}
\end{lstlisting}

\ui{Apply}, then \ui{Typeset}.

\subsection{Worked example: tightening a paragraph}

Select a paragraph you are unhappy with and run \ui{Proofread Section}. A
typical diff:

\begin{lstlisting}
- The experiments that we ran showed good results across all
- three of the datasets that we evaluated on.
+ Our experiments show strong results across all three datasets.
\end{lstlisting}

Because the skill also receives your labels and citation keys, a suggestion to
cite something will name a key you actually have.

Read every diff before applying. The assistant is a fast, well-informed
copy-editor, not an authority on your argument --- and on a rewrite skill it
will happily reword a sentence whose awkwardness was deliberate.

\subsection{Copy Prompt}

\ui{Copy Prompt} puts the exact prompt --- context included --- on your
clipboard, for pasting into another tool. It is the same text \ui{Run} would
send, so what you get elsewhere matches what you would have got here.

\subsection{What it is actually good at}
\label{sec:ai-measured}

Every section of this guide was run through \ui{Proofread Section}. The results
are worth stating plainly, because they describe the tool better than any
feature list.

\paragraph{It does not break your \LaTeX.}
Across all six sections, every \verb|\begin| matched its \verb|\end|, and every
\verb|\label| and \verb|\ref| came back intact. That is the property that makes
\ui{Apply} safe to press.

\paragraph{On polished prose, it finds almost nothing.}
Six sections produced exactly one correction worth taking --- a missing comma
between two independent clauses:

\begin{lstlisting}
- It does not upload your documents anywhere and it does not
+ It does not upload your documents anywhere, and it does not
\end{lstlisting}

Everything else was house style rather than error: serial commas, added in some
lists and not others; and removing the spaces around em-dashes, which this
document uses consistently, so accepting it would have made one section
disagree with the other five. Applied wholesale, those edits would have left
the guide \emph{less} consistent than it started.

\paragraph{Watch for reflowing.}
The largest ``change'' was not a change at all. On one section the model
returned text that was word-for-word identical to the input, rewrapped onto
different lines --- forty-eight changed lines, no difference in the PDF. The
skills now instruct the model to preserve line breaks, which reduced that to
zero, but it is the failure mode to recognise: a diff that is enormous and
entirely uniform is usually reflowing, not editing.

\paragraph{So what is it for?}
Mechanics, on text that is still rough. A first draft, notes pasted in from
elsewhere, a paragraph written at midnight --- there it earns its place
immediately. On prose you have already revised twice, expect it to propose
preferences and expect to decline most of them. That is not a failure of the
model; it is what running a copy-editor over clean copy looks like.

Which is why every rewrite skill ends in a diff and a button. Judge each
changed line on its own: accepting a whole rewrite because most of it is
harmless is how a document loses its voice one paragraph at a time.

\subsection{When it goes wrong}

If a request fails, the error appears in its own red panel and the result area
stays empty --- an error never masquerades as output you might paste into your
document. \ui{Cancel} stops a run in progress and keeps whatever streamed in so
far, in case the useful part arrived early.
